\documentclass[11pt]{article}

\usepackage[margin=1in]{geometry}
\usepackage{amsmath,amssymb,amsthm,mathtools}
\usepackage{graphicx}
\usepackage[hidelinks]{hyperref}
\usepackage{enumitem}

\setlength{\parindent}{0pt}
\setlength{\parskip}{0.55em}
\setlist[itemize]{leftmargin=1.5em}

\newtheorem{theorem}{Theorem}
\newtheorem{lemma}{Lemma}
\newtheorem{corollary}{Corollary}
\newtheorem{remark}{Remark}

\newcommand{\BH}{\mathcal B(\mathcal H)}
\newcommand{\BB}{\mathfrak B}
\newcommand{\Spw}{\mathcal S^p_w}
\newcommand{\norm}[1]{\left\lVert #1\right\rVert}
\newcommand{\snorm}[2]{\left\lVert #1\right\rVert_{#2}}
\newcommand{\rad}{r}

\begin{document}

\begin{center}
{\Large Full Solution Packet}\\[0.3em]
{\large Operator-valued Schur classes satisfy Sun's Wiener lemma}\\[0.45em]
\textbf{Status:} candidate full result, likely valid\\
\textbf{Target:} Remark 3.11(2) of arXiv:2407.16416
\end{center}

\section*{Source and exact future-work problem}

The source is K\"ohldorfer--Balazs,
\emph{Wiener pairs of Banach algebras of operator-valued matrices},
arXiv:2407.16416, published in the \emph{Journal of Mathematical Analysis
and Applications} 549 (2025), 129525 \cite{KB}.  On page 18, Remark
3.11(2) says:

\begin{center}
\includegraphics[width=\textwidth]{figures/open_problem_crop.png}
\end{center}

In text: Sun proved inverse-closedness of the scalar-valued versions of
$\mathcal S^p_\nu(X)$ in $\mathcal B(\ell^2(X;\mathbb C))$.  The authors
expect analogous results for $\BH$-valued matrices and leave the detailed
investigation to future research.

We prove the operator-valued analogue of Sun's theorem under his original
admissibility hypothesis: the weight is $(p,2)$-admissible.  The local
dimension of $\mathcal H$ is arbitrary, including infinite.

\section*{Definitions and hypotheses}

Let $(X,\rho)$ be a discrete pseudometric space whose counting measure has
polynomial growth, and let $\mathcal H$ be a Hilbert space.  For
$1\le p\le\infty$ and a symmetric weight $w$ on $X\times X$, define
$\Spw(X;\BH)$ to consist of block matrices $A=(A_{xy})_{x,y\in X}$ with
$A_{xy}\in\BH$ and
\[
 \snorm{A}{p,w}:=
 \max\left\{
   \sup_x\norm{\bigl(\norm{A_{xy}}w(x,y)\bigr)_y}_{\ell^p},
   \sup_y\norm{\bigl(\norm{A_{xy}}w(x,y)\bigr)_x}_{\ell^p}
 \right\}<\infty.
\]
For $p=\infty$ this has the usual supremum interpretation.  This is the
operator-valued space $\mathcal S^p_\nu(X)$ in \cite{KB} when
$w(x,y)=\nu(x-y)$.

We use Sun's definition of a $(p,r)$-admissible weight
\cite[Definition in Section 2.2]{Sun}.  Thus there are a companion weight
$v$, constants $D>0$ and $0<\theta<1$, and, writing $p'$ and $r'$ for
conjugate exponents, the following properties:
\begin{align}
 w(x,y)&\le D\bigl(w(x,z)v(z,y)+v(x,z)w(z,y)\bigr),
                                                        \label{eq:weight-product}\\
 \sup_x\norm{(v/w)(x,\cdot)}_{p'}
 +\sup_y\norm{(v/w)(\cdot,y)}_{p'}&<\infty,             \label{eq:embedding-weight}\\
 \inf_{\tau>0}\bigl(a_{r'}(\tau)+b_{p'}(\tau)t\bigr)
 &\le D t^\theta\qquad(t\ge1).                         \label{eq:interpolation-weight}
\end{align}
Here $a_{r'}(\tau)$ is the uniform $\ell^{r'}$ size of $v$ on the
$\rho$-ball of radius $\tau$, in rows and columns, while $b_{p'}(\tau)$ is
the uniform $\ell^{p'}$ size of $v/w$ outside that ball, again in rows and
columns.  These are exactly Sun's formulas (2.8)--(2.12); we have recalled
their content to make clear where the proof uses the hypothesis.

Sun notes immediately after the definition that $(p,r)$-admissibility
implies $(p,\widetilde r)$-admissibility whenever $r\le\widetilde r$.  The
single crucial move below is therefore legitimate:
\begin{equation}
        (p,2)\text{-admissible}\quad\Longrightarrow\quad
        (p,\infty)\text{-admissible}.
                                                        \label{eq:move-r}
\end{equation}

\section*{Proof intuition}

Sun's scalar proof controls a local row/column norm of a matrix power by an
ambient operator norm.  At exponent $r=2$, its delta-vector argument is false
for operator blocks: many norm-one blocks in one row may act on mutually
orthogonal directions, so their scalar $\ell^2$ norm can be large although
the block-row operator has norm one.

The admissibility hypothesis contains the repair.  Run the same interpolation
with auxiliary exponent $r=\infty$.  Every individual block is a coordinate
compression of the global operator, hence
\[
        \sup_{x,y}\norm{T_{xy}}
        \le \norm{T}_{\mathcal B(\ell^2(X;\mathcal H))}.
\]
This estimate is valid for arbitrary block size.  The resulting
interpolation inequality makes $\Spw$ a differential subalgebra of the
ambient operator algebra.  A dyadic spectral-radius argument and Hulanicki's
lemma then give inverse-closedness.

\section*{The differential estimate}

Set
\[
        \BB=\mathcal B(\ell^2(X;\mathcal H)).
\]

\begin{lemma}[Operator-valued Sun estimate]\label{lem:differential}
Assume that $w$ is $(p,2)$-admissible.  Then $\Spw(X;\BH)$ is a unital
Banach $*$-algebra continuously embedded in $\BB$.  Moreover, for some
constants $C>0$ and $0<\theta<1$ depending only on the admissibility data,
\begin{equation}
 \snorm{T^2}{p,w}
 \le C\,\snorm{T}{p,w}^{,1+\theta}
          \norm{T}_{\BB}^{,1-\theta}
 \qquad(T\in\Spw).
                                                        \label{eq:differential}
\end{equation}
\end{lemma}

\begin{proof}
Completeness of the operator-valued space is Proposition 3.2 of \cite{KB}.
We prove the algebra, embedding, and differential estimates.

Matrix multiplication and \eqref{eq:weight-product}, followed by the standard
weighted Jensen--Schur estimate for composing an $\ell^p$ row/column matrix
with an $\ell^1$ row/column matrix, give the compensated product estimate
\begin{equation}
 \snorm{RS}{p,w}
 \le C_0\bigl(
       \snorm{R}{p,w}\snorm{S}{1,v}
       +\snorm{R}{1,v}\snorm{S}{p,w}
       \bigr).                                          \label{eq:product}
\end{equation}
This is Sun's estimate (3.3).  It remains valid for operator blocks because
the only entrywise input is
\[
 \norm{(RS)_{xy}}
 \le\sum_z\norm{R_{xz}}\norm{S_{zy}}.
\]
No commutativity is used.

By H\"older's inequality and \eqref{eq:embedding-weight},
\begin{equation}
        \snorm{T}{1,v}\le C_1\snorm{T}{p,w}.             \label{eq:embedding}
\end{equation}
Since $v\ge1$, the operator-valued Schur test then yields
\begin{equation}
        \norm{T}_{\BB}\le C_1\snorm{T}{p,w}.            \label{eq:ambient-embedding}
\end{equation}
Equations \eqref{eq:product} and \eqref{eq:embedding} prove the Banach
algebra property.  Symmetry of $w$ gives
$\snorm{T^*}{p,w}=\snorm{T}{p,w}$, and the diagonal identity belongs to the
algebra because $X$ has counting measure and $w(x,x)$ is uniformly bounded.

We now use $(p,\infty)$-admissibility from \eqref{eq:move-r}.  Split each
row and column of $T$ into $\rho\le\tau$ and $\rho>\tau$.  On the near part,
use the block-compression estimate
\begin{equation}
 \norm{T_{xy}}
 =\norm{P_xTP_y}\le\norm{T}_{\BB},                      \label{eq:block-compression}
\end{equation}
where $P_x$ is the coordinate projection onto the $x$-copy of $\mathcal H$.
On the far part, use H\"older's inequality with $w$ and $v/w$.  Taking the
row and column suprema gives
\begin{equation}
 \snorm{T}{1,v}
 \le a_1(\tau)\norm{T}_{\BB}
       +b_{p'}(\tau)\snorm{T}{p,w}.                     \label{eq:split}
\end{equation}

Let $K=\max\{1,C_1\}$ in \eqref{eq:ambient-embedding} and put
$t=K\snorm{T}{p,w}/\norm{T}_{\BB}$ when $T\ne0$.  Then $t\ge1$.
Take the infimum in \eqref{eq:split} and apply
\eqref{eq:interpolation-weight}; enlarging constants to absorb $K$ gives
\begin{equation}
 \snorm{T}{1,v}
 \le C_2\norm{T}_{\BB}^{,1-\theta}
             \snorm{T}{p,w}^{\,\theta}.                 \label{eq:interpolation}
\end{equation}
The zero operator is trivial.  Finally apply \eqref{eq:product} with
$R=S=T$ and then \eqref{eq:interpolation}.  This yields
\eqref{eq:differential}.
\end{proof}

\section*{Main result}

\begin{theorem}[Operator-valued Wiener lemma]\label{thm:main}
Let $1\le p\le\infty$, let $(X,\rho)$ be a discrete pseudometric space with
counting measure of polynomial growth, and let $w$ be a $(p,2)$-admissible
weight in the sense of Sun.  For every Hilbert space $\mathcal H$,
\[
 \mathcal S^p_w(X;\BH)
 \subseteq \mathcal B(\ell^2(X;\mathcal H))
\]
is inverse-closed.  Equivalently, if $A\in\mathcal S^p_w(X;\BH)$ is
invertible on $\ell^2(X;\mathcal H)$, then
\[
        A^{-1}\in\mathcal S^p_w(X;\BH).
\]
The algebra $\mathcal S^p_w(X;\BH)$ is also symmetric.
\end{theorem}

\begin{proof}
Write $\mathcal A=\Spw(X;\BH)$.  The continuous unital inclusion
$\mathcal A\hookrightarrow\BB$ always gives
\begin{equation}
        \rad_{\BB}(T)\le\rad_{\mathcal A}(T).            \label{eq:easy-radius}
\end{equation}
We prove the reverse inequality.

Apply Lemma \ref{lem:differential} to $T^{2^n}$.  With
\[
 A_n=\snorm{T^{2^n}}{p,w}^{1/2^n},
 \qquad
 B_n=\norm{T^{2^n}}_{\BB}^{1/2^n},
 \qquad
 \alpha=\frac{1+\theta}{2}<1,
\]
we obtain
\begin{equation}
 A_{n+1}
 \le C^{1/2^{n+1}} A_n^\alpha B_n^{1-\alpha}.           \label{eq:dyadic}
\end{equation}
The sequence $(A_n)$ is bounded (equivalently, replace the algebra norm by
an equivalent submultiplicative norm), and Gelfand's formula gives
$B_n\to\rad_{\BB}(T)$.  If $\lambda>\rad_{\BB}(T)$, then $B_n\le\lambda$
for all large $n$.  Taking the limsup in \eqref{eq:dyadic} yields
\[
 L\le L^\alpha\lambda^{1-\alpha},
 \qquad L:=\limsup_n A_n.
\]
Thus either $L=0$ or $L\le\lambda$.  Letting
$\lambda\downarrow\rad_{\BB}(T)$ gives
\[
        \rad_{\mathcal A}(T)\le\rad_{\BB}(T).
\]
Together with \eqref{eq:easy-radius}, the spectral radii agree (in
particular, they agree for every self-adjoint $T\in\mathcal A$).

Hulanicki's lemma for a unital Banach $*$-subalgebra of a $C^*$-algebra now
implies that $\mathcal A$ is inverse-closed in $\BB$ and is symmetric; see
\cite[Proposition 2.1]{KB}.  This proves the theorem.
\end{proof}

\begin{corollary}[Polynomial off-diagonal weights]\label{cor:polynomial}
Let
\[
 d_X=\inf\{d:\ \#B_\rho(x,R)\le C R^d
                    \text{ uniformly in }x,\ R\ge1\}.
\]
For $s>d_X(1-1/p)$, put $w_s(x,y)=(1+\rho(x,y))^s$.  Then
$\mathcal S^p_{w_s}(X;\BH)$ is inverse-closed in
$\mathcal B(\ell^2(X;\mathcal H))$.
In particular, for a relatively separated $X\subset\mathbb R^d$, the
sufficient condition $s>d(1-1/p)$ gives the operator-valued extension
requested in Remark 3.11(2) of \cite{KB}.
\end{corollary}

\begin{proof}
Sun's Example A.2 proves that $w_s$ is $(p,r)$-admissible under the displayed
condition, hence in particular $(p,2)$-admissible.  Apply Theorem
\ref{thm:main}.
\end{proof}

The same argument applies to the subexponential weights in Sun's Example
A.3 (under its metric-space hypothesis).

\section*{Verification notes}

The proof was audited at the four places where scalar-to-block extensions
can fail.

\begin{itemize}
\item The entrywise product estimate uses only
      $\norm{RS}\le\norm R\norm S$ in $\BH$, so noncommutativity is harmless.
\item The scalar $r=2$ delta-vector estimate was \emph{not} imported.  It is
      replaced by $r=\infty$ and the exact compression inequality
      \eqref{eq:block-compression}.
\item Sun explicitly records the monotonic implication
      $(p,2)$-admissible $\Rightarrow(p,\infty)$-admissible, so the changed
      auxiliary exponent does not strengthen the theorem's hypothesis.
\item The passage from the differential estimate to inverse-closedness is
      the dyadic spectral-radius argument above followed by the same
      Hulanicki lemma used in the source paper.
\end{itemize}

No computational or finite-dimensional test is used as evidence.  The proof
is uniform in $\dim\mathcal H$.

\section*{Novelty check and limitations}

Before promotion, the run indexes were searched for arXiv:2407.16416,
``Wiener pairs'', ``operator-valued Schur'', ``inverse-closed'', and
``$\mathcal S^p_\nu$''.  No duplicate solution or attempt was found.

A bounded web search on 2026-07-19 used the exact source title and phrases
combining ``operator-valued'', ``Schur class'', ``Wiener lemma'', and
``inverse-closed''.  It found the 2025 journal version of the source and the
authors' arXiv:2501.09603 follow-up on the operator-valued Jaffard algebra,
but no paper stating the $\mathcal S^p_w$ extension proved here.  The search
also checked Sun's original 2007 paper and author-hosted PDF.  Novelty is
therefore plausible, not certified.

The theorem uses exactly the $(p,2)$-admissibility hypothesis of Sun's scalar
result.  It does not claim inverse-closedness for arbitrary weights for which
$\mathcal S^p_w$ is merely defined.  This matches the natural precise reading
of the source's requested ``analogous results.''

\section*{Human review recommendation}

Send the packet to an expert in inverse-closed matrix algebras.  The most
important review point is the interpolation step
\eqref{eq:split}--\eqref{eq:interpolation}: verify the convention for Sun's
$a_{r'}$ and $b_{p'}$ and the monotonic passage from $r=2$ to $r=\infty$.
After that, the block-compression and spectral-radius arguments are
dimension-free and elementary.

\begin{thebibliography}{9}

\bibitem{KB}
Lukas K\"ohldorfer and Peter Balazs,
\emph{Wiener pairs of Banach algebras of operator-valued matrices},
J. Math. Anal. Appl. 549 (2025), 129525;
arXiv:2407.16416.

\bibitem{Sun}
Qiyu Sun,
\emph{Wiener's lemma for infinite matrices},
Trans. Amer. Math. Soc. 359 (2007), 3099--3123.

\end{thebibliography}

\end{document}
